\section{Discussion}
\label{sec:discussion}

\paragraph{Choosing a model} The leading API models are not significantly different on this benchmark, so price can decide among them: \ThirdName{} gives \ScoreGptSixOneSol{} for US\$\ThirdUSD{} per 100 items and gemini-3.7-flash \ScoreGeminiThreeSevenFlash{} for US\$\UsdGeminiThreeSevenFlash{}, against \TopScore{} for US\$\TopUSD{}; \CheapName{} (\CheapScore{} for US\$\CheapUSD{}) is the economical choice for high-volume use. The cheapest good models are not the ones with \emph{mini} or \emph{flash} in their names, and the most expensive model was not a flagship: measured cost per item is a better guide than tier labels.

\paragraph{References need a checker, whatever the model} \NoneApiNonzero{} of the \NApi{} API models listed references that no database knows, and many listed real works with a wrong DOI, year or author (Fig.~\ref{fig:references}). A reference produced by a model should be verified before it enters a thesis; a lookup in Crossref, with a second database for the misses, costs next to nothing compared with a fabricated citation in an examined document. Thesis tools should treat references as a retrieval problem, not a generation problem.

\paragraph{Local models} Local models are the option when data may not leave the institution, and for the tasks where they come close to API models: rewriting text into formal Indonesian and answering methodology questions. They were not cheaper per item at rental prices than the cheapest API models of much higher quality, and they should not be trusted with references. The model tuned for Indonesian did not outperform a newer general model of its size, which suggests that for these tasks a recent general model is a better starting point than language specialisation alone (one pair of models with different quantisation; not a controlled comparison).

\paragraph{Benchmark design} Two tasks are near ceiling for API models (methodology, formal Indonesian) while still separating local models; harder variants---open questions, longer drafts with subtler non-standard forms---would extend their range. Judging was the largest cost item. Scoring the prose tasks by the deterministic checks alone barely changes the overall ranking (Spearman $\rho=\JudgeAutoRho{}$, largest rank shift \JudgeAutoShift{}), so judging a stratified sample of answers may suffice; we have not tested this.

\paragraph{Limitations} Each model answered each item once, at low reasoning effort, so run-to-run variation is not measured; the \NPremium{} premium-tier models (estimated at about US\$50) and a high-effort check of the best models are not in this version. AI units are the vendor's billing measure at one point in time, hosted models change without notice, and the served-model check relies on the SDK's report. The structure references and the extraction excerpts come from the theses of one university; extraction is scored against silver labels, which for a candidate come from the panel members of other vendors and therefore differ between vendors; the extraction items are private, but without them the ranking is almost the same (Spearman $\rho=\PubRho{}$, largest rank shift \PubMaxShift{}). Neither the LLM judges nor the item keys have yet been checked by people, and the items were drafted with LLM assistance, which may favour LLM phrasing; deterministic scoring limits this risk but does not remove it. The grounded checks test citation form, coverage and numbers, not whether each cited claim is supported, which is left to the judges. A real work that none of the four databases knows counts as not found. Local models ran 4-bit quantised in one serving stack, their energy was measured on a shared server, and their cost depends on the concurrency used.
